\documentclass[11pt,reqno]{amsart}
%DIF LATEXDIFF DIFFERENCE FILE
%DIF DEL homogeneous_arveson_douglas_announcement_v2.tex   Tue Nov 24 18:40:31 2026
%DIF ADD homogeneous_arveson_douglas_announcement_v3.tex   Tue Nov 24 18:40:31 2026
\usepackage[T1]{fontenc}
\usepackage[utf8]{inputenc}
\usepackage{lmodern}
\usepackage[letterpaper,margin=1.02in]{geometry}
\usepackage{amsmath,amssymb,amsfonts,mathtools,mathrsfs}
\usepackage{microtype}
\usepackage[dvipsnames]{xcolor}
\usepackage{hyperref}
\hypersetup{hidelinks,pdftitle={The homogeneous Arveson--Douglas conjecture: a research announcement},
 pdfsubject={Vector-valued homogeneous Arveson--Douglas theorem and a three-step projective lifting outline}}
\newtheorem{theorem}{Theorem}
\renewcommand{\thetheorem}{\Alph{theorem}}
\theoremstyle{remark}
\newtheorem*{remark}{Remark}
\newcommand{\C}{\mathbb C}
\newcommand{\N}{\mathbb N}
\newcommand{\PP}{\mathbb P}
\newcommand{\cO}{\mathcal O}
\newcommand{\cI}{\mathscr I}
\newcommand{\cK}{\mathscr K}
\newcommand{\cS}{\mathcal S}
\newcommand{\eps}{\varepsilon}
\newcommand{\dbar}{\bar\partial}
\newcommand{\FS}{\mathrm{FS}}
\newcommand{\DA}{\mathrm{DA}}
\DeclareMathOperator{\im}{im}
\DeclareMathOperator{\Quot}{Quot}
\DeclareMathOperator{\Gr}{Gr}
\DeclareMathOperator{\Supp}{Supp}
\DeclareMathOperator{\Id}{Id}
\allowdisplaybreaks[2]
\emergencystretch=1.2em
\title[The homogeneous Arveson--Douglas conjecture]{The homogeneous Arveson--Douglas conjecture:\\a research announcement}
\date{23 November 2026}
\thanks{Research draft for independent mathematical assessment. The statements
 announced here summarize the proposed proof in the companion manuscript
%DIF 39c39
%DIF <  \cite{DingFangTang}; this announcement is not a substitute for its detailed proofs.}
%DIF -------
 \textcolor{red}{\cite{DingFangTang}} \textcolor{blue}{\cite{Companion}}; this announcement is not a substitute for its detailed proofs.} %DIF > 
%DIF -------
\subjclass[2020]{Primary 47A13; Secondary 32W05, 13D02, 47B10}
%DIF < \keywords{Arveson--Douglas conjecture, homogeneous submodules, Drury--Arveson space, projective geometry, holomorphic division, essential normality}
%DIF -------
\keywords{Arveson--Douglas conjecture, homogeneous submodules, Drury--Arveson space, %DIF > 
 projective geometry, holomorphic division, essential normality} %DIF > 
%DIF PREAMBLE EXTENSION ADDED BY LATEXDIFF
%DIF CFONT PREAMBLE %DIF PREAMBLE
\RequirePackage{color}\definecolor{RED}{rgb}{1,0,0}\definecolor{BLUE}{rgb}{0,0,1} %DIF PREAMBLE
\DeclareOldFontCommand{\sf}{\normalfont\sffamily}{\mathsf} %DIF PREAMBLE
\providecommand{\DIFaddtex}[1]{{\protect\color{blue} #1}} %DIF PREAMBLE
\providecommand{\DIFdeltex}[1]{{\protect\color{red} #1}} %DIF PREAMBLE
%DIF SAFE PREAMBLE %DIF PREAMBLE
\providecommand{\DIFaddbegin}{} %DIF PREAMBLE
\providecommand{\DIFaddend}{} %DIF PREAMBLE
\providecommand{\DIFdelbegin}{} %DIF PREAMBLE
\providecommand{\DIFdelend}{} %DIF PREAMBLE
\providecommand{\DIFmodbegin}{} %DIF PREAMBLE
\providecommand{\DIFmodend}{} %DIF PREAMBLE
%DIF FLOATSAFE PREAMBLE %DIF PREAMBLE
\providecommand{\DIFaddFL}[1]{\DIFadd{#1}} %DIF PREAMBLE
\providecommand{\DIFdelFL}[1]{\DIFdel{#1}} %DIF PREAMBLE
\providecommand{\DIFaddbeginFL}{} %DIF PREAMBLE
\providecommand{\DIFaddendFL}{} %DIF PREAMBLE
\providecommand{\DIFdelbeginFL}{} %DIF PREAMBLE
\providecommand{\DIFdelendFL}{} %DIF PREAMBLE
%DIF HYPERREF PREAMBLE %DIF PREAMBLE
\providecommand{\DIFadd}[1]{\texorpdfstring{\DIFaddtex{#1}}{#1}} %DIF PREAMBLE
\providecommand{\DIFdel}[1]{\texorpdfstring{\DIFdeltex{#1}}{}} %DIF PREAMBLE
%DIF COLORLISTINGS PREAMBLE %DIF PREAMBLE
\RequirePackage{listings} %DIF PREAMBLE
\RequirePackage{color} %DIF PREAMBLE
\lstdefinelanguage{DIFcode}{ %DIF PREAMBLE
%DIF DIFCODE_CFONT %DIF PREAMBLE
  moredelim=[il][\color{red}\scriptsize]{\%DIF\ <\ }, %DIF PREAMBLE
  moredelim=[il][\color{blue}\sffamily]{\%DIF\ >\ } %DIF PREAMBLE
} %DIF PREAMBLE
\lstdefinestyle{DIFverbatimstyle}{ %DIF PREAMBLE
	language=DIFcode, %DIF PREAMBLE
	basicstyle=\ttfamily, %DIF PREAMBLE
	columns=fullflexible, %DIF PREAMBLE
	keepspaces=true %DIF PREAMBLE
} %DIF PREAMBLE
\lstnewenvironment{DIFverbatim}{\lstset{style=DIFverbatimstyle}}{} %DIF PREAMBLE
\lstnewenvironment{DIFverbatim*}{\lstset{style=DIFverbatimstyle,showspaces=true}}{} %DIF PREAMBLE
%DIF END PREAMBLE EXTENSION ADDED BY LATEXDIFF

\begin{document}
\begin{abstract}
We announce a projective lifting approach to the homogeneous
Arveson--Douglas conjecture for homogeneous submodules of the vector-valued
Drury--Arveson space $H_m^2\otimes\C^r$, with $r<\infty$.
The principal estimate bounds the degree-$n$ coordinate defect by
$C/(n+1)$. For an infinite-dimensional algebraic quotient of support
dimension $\delta$, it gives quotient commutators in the weak Schatten
class $\cS_{\delta,\infty}$. The outline has three steps: realize the
graded scalar problem on projective space with its exact norm;
construct a uniform holomorphic lift by division, compatible metrics,
and a finite $\bar\partial$ correction; and return the estimate to the
Drury--Arveson shift, including finite coefficient multiplicity through
auxiliary variables. Singular values are counted in the original
vector quotient. The intermediate constructions are stated as
theorems, and a historical account identifies their division,
resolution, and Dolbeault antecedents.
\end{abstract}
\maketitle

\section{Historical background}

The Arveson--Douglas problem asks how polynomial relations affect the
essential normality of the multivariable shift. Arveson's 1998 work
established the Drury--Arveson space as a model for commuting row
contractions and proved Schatten estimates for the full shift, together
with its sphere Toeplitz extension \cite{Arveson1998}. Compression to a
polynomial quotient introduces a further term: the part of coordinate
multiplication discarded by the quotient projection. Controlling this
term is the analytic content of the problem.

For submodules generated by vector monomials, Arveson subsequently
proved projection-commutator estimates in $\cS_p$ for $p>2m$ and
quotient-commutator estimates for $p>m$, where $m$ is the number of
variables. His Conjecture~A asks for compactness of the projection
commutators for arbitrary homogeneous vector generators
\cite{Arveson2005}. Douglas connected essential normality with index
theory and $K$-homology, and proposed that the dimension of the quotient's
support should determine its summability threshold \cite{DouglasIndex}.
His Conjecture~9 is stated for ball Bergman quotients and predicts
$p>\dim(Z\cap\mathbb B^m)$. The homogeneous Drury--Arveson formulation
has the same dimension principle.

Guo and Kai Wang established important homogeneous cases, including
principal submodules in arbitrary ambient dimension and the scalar
problem in two and three variables \cite{GuoWang}. Shalit made
norm-controlled polynomial division an explicit organizing principle
\cite{ShalitDivision}. A controlled decomposition
$F=\sum_\nu a_\nu f_\nu$, followed by differentiation and projection
perpendicular to the ideal, eliminates the terms
$\sum_\nu f_\nu\partial a_\nu$. The differentiated generators then
give an off-diagonal estimate of order $n^{-1/2}$, and the Hilbert
function supplies the dimension-dependent Schatten threshold.
Eschmeier's related operator-inequality criterion
\cite{EschmeierCriterion}, in the projected-derivative form recorded
in \cite[Section~6, equation~(6.1)]{ShalitDivision}, identifies the
estimate that a division argument must produce.
Shalit's scalar encoding of vector polynomials by expressions linear
in auxiliary variables also supplies the antecedent for the passage
to finite coefficient multiplicity
\cite[Section~5, equation~(5.1)]{ShalitDivision}.

Engli\v{s} and Eschmeier proved the geometric conjecture for homogeneous
varieties smooth away from the origin \cite{EnglisEschmeier}. Their
submodule consists of functions vanishing on the variety, corresponding
to radical ideals. Their method uses restriction to the variety and
the generalized Toeplitz calculus of Boutet de Monvel and Guillemin,
within a scale of holomorphic Sobolev spaces containing the
Drury--Arveson space. For a variety of dimension $\delta$, their
homogeneous proof gives $\cS_p$ membership for $p>\delta$ and explicitly
records the weak endpoint $\cS_{\delta,\infty}$. They also treat smooth
components of different dimensions meeting only at the origin. Thus
both the geometric interpretation and the weak Schatten endpoint have
substantial precedents.

Douglas, Tang, and Yu developed an $L^2$ Dolbeault and index-theoretic
approach in their analytic Grothendieck--Riemann--Roch theorem. In their complete-intersection setting, the
variety meets the sphere transversely, its boundary is smooth, and
its interior singularities may be isolated. Their Section~3 uses $\dbar$-Neumann
analysis.
This analytical framework produces on the regular locus to construct a relative Dolbeault
class whose
boundary is the CR fundamental class. For $I=(p_1,\ldots,p_M)$,
$M\le m-2$, with the generator Jacobian of rank $M$ along
$Z_I\cap\partial\mathbb B^m$ and transverse sphere intersection,
their complete-intersection setting has smooth boundary and possibly
isolated interior singularities. A bounded restriction and extension
theorem connects the weighted Bergman module on the variety. Section~4 connects the variety's weighted Bergman space to
the ball quotient by restriction and a Beatrous-type integral
extension. They obtain essential normality of the ideal and
quotient and identify the quotient
Toeplitz class with the boundary's $\mathrm{spin}^c$ fundamental class
of the boundary \cite[Theorems~1.2, 1.4 and~1.5]{DouglasTangYu}.
Homogeneity is not required. Both their $\dbar$ analysis and their
restriction and
extension methods are substantive antecedentsof the
present argument. This nonhomogeneous geometric result supplies both Dolbeault and
extension antecedents; its integral extension construction and the
finite correction used below are distinct analytic procedures.

Douglas, Jabbari, Tang, and Yu developed the resolution viewpoint
outlined by Douglas \cite{DouglasIndex} into a construction for
monomial ideals \cite{DouglasJabbariTangYu}. They give finite exact
complexes of essentially normal Hilbert modules with bounded module
maps, and express the quotient Toeplitz class as an alternating sum
of the classes of the terms. The construction retains normal-jet
information and applies also to the Drury--Arveson space. Thus the
use of resolutions to connect essential normality, nonreduced
structure, and index theory is an established part of the subject.

Ding \cite{Ding26} extended Engli\v{s} and Eschmeier's results  to strongly pseudoconvex  manifold with the Property (S) by Hironaka’s embedded desingularization theorem, and considered the  Schatten-class problem of the antisymmetric sum of the compression of the multiplication operators to the quotient  holomorphic Sobolev spaces.
The antisymmetric sum was induced by Holton and Howe \cite{HeltonHowe75}, and plays a key role in noncommutative geometry. The division and correction steps also belong to a longer history
in complex analysis. Parameter-dependent division and quantitative
linear splittings were developed by Galligo and Hauser--M\"uller
\cite{Galligo,HauserMuller}. Holomorphic division by $L^2$ methods
goes back to H\"ormander, Kelleher--Taylor, and Skoda
\cite{HormanderGenerators,KelleherTaylor,Skoda1972,Skoda1978}.
Andersson--Wulcan's treatment of graded resolutions on projective
space is a particularly close antecedent of the descending
$\dbar$ construction \cite[Section~6]{AnderssonWulcan}.
The discussion of Theorems~\ref{thm:division}--\ref{thm:lift} below
specifies what is inherited from these works and what quantitative
assertions are needed here.

The companion manuscript \cite{DingFangTang, Companion} develops these division,
resolution, and Dolbeault methods into a uniform projective lifting
construction for homogeneous ideal sheaves. It works on smooth
ambient projective space while retaining nilpotent and embedded
structure. Complete spaces of sections control degeneration; sampled
coefficient metrics place division, differentiated generators, and
curvature in compatible norms. Positive twisting and a finite
correction through the syzygies then give a global lift bounded in
the original scalar norm. An auxiliary-variable construction transfers
the resulting estimate to vector-valued modules. The additional claim
is the uniform analytic control that makes this passage effective;
the operator reduction and the role of the Hilbert function retain
their established precedents.

\section{The vector-valued theorem and the three-step outline}

Let $m,r\ge1$, let $R=\C[z_1,\ldots,z_m]$, and give the Drury--Arveson space $H_m^2$
its standard normalization
\[
 \langle z^\alpha,z^\beta\rangle=0\quad(\alpha\ne\beta),
 \qquad \|z^\alpha\|^2=\frac{\alpha!}{|\alpha|!}.
\]
Fix a Hermitian coefficient space $\C^r$, placed in degree zero.
For a homogeneous submodule $M\subset R\otimes\C^r$, set
\[
 \mathcal M=\overline M\subset H_m^2\otimes\C^r,\qquad
 Q=\mathcal M^\perp,\qquad Q_n=(R_n\otimes\C^r)\ominus M_n.
\]
Write $Z_j=M_{z_j}\otimes\Id_r$, $S_j=P_QZ_j|_Q$, and
$V_j=Z_j|_{\mathcal M}$. Introduce the off-diagonal maps and their
positive defect
\begin{equation}\label{eq:defect}
 B_{j,n}=P_{M_{n+1}}Z_j|_{Q_n},
 \qquad A_n=\sum_{j=1}^m B_{j,n}^*B_{j,n}.
\end{equation}
For $p>0$, the weak Schatten class $\cS_{p,\infty}$ consists of compact
operators with singular values $s_k(T)=O(k^{-1/p})$.

\begin{theorem}[Homogeneous essential normality in finite multiplicity]\label{thm:main}
For every homogeneous submodule $M\subset R\otimes\C^r$ there are
constants $C_M,C_M'<\infty$ such that, for all $n\ge0$,
\begin{equation}\label{eq:main}
 0\le A_n\le\frac{C_M}{n+1}\Id_{Q_n},
 \qquad
 \|[S_i^*,S_j]|_{Q_n}\|\le\frac{C_M'}{n+1}.
\end{equation}
If $T=(R\otimes\C^r)/M$ is infinite dimensional and
$\delta=\dim\Supp_R T$, then
\begin{equation}\label{eq:weak}
 [S_i^*,S_j]\in\cS_{\delta,\infty}
 \subset\bigcap_{p>\delta}\cS_p.
\end{equation}
The column $B=(P_{\mathcal M}Z_j|_Q)_{j=1}^m:Q\to\mathcal M^{\oplus m}$
and the projection commutators $[P_{\mathcal M},Z_j]$ belong to
$\cS_{2\delta,\infty}$. The restriction commutators satisfy
\[
 [V_i^*,V_j]\in\cS_{m,\infty}
 \subset\bigcap_{p>m}\cS_p.
\]
If $T$ is finite dimensional, the quotient commutators and off-diagonal
operators have finite rank; the restriction assertion still holds.
\end{theorem}

The support dimension is the complex dimension of
$\Supp_R T\subset\C^m$. The constants depend on the fixed module and
its original Hermitian geometry. The statement uses the standard
grading and includes every closed graded submodule of
$H_m^2\otimes\C^r$: its homogeneous polynomial part is finitely
generated and dense. For $r=1$ it includes arbitrary homogeneous
ideals, with singular or nonreduced structure.

The proof separates the size of an operator block from its multiplicity.
The exact identities
\begin{align}
 [S_i^*,S_j]|_{Q_n}
 &=P_Q[Z_i^*,Z_j]|_{Q_n}-B_{i,n}^*B_{j,n},\label{eq:compression}\\
 [Z_i^*,Z_j]|_{R_n\otimes\C^r}
 &=\frac{\delta_{ij}\Id-Z_jZ_i^*}{n+1}\label{eq:ambient}
\end{align}
reduce the analytic task to the first estimate in \eqref{eq:main}.
The Hilbert function then gives the exponent $\delta$.

\subsection*{The proof in three steps}
\begin{enumerate}
\item \textbf{Projective realization.}
Homogeneity identifies polynomial degrees with powers of the
hyperplane bundle. Saturation and an exact norm identity transfer the
scalar problem to an ideal sheaf on smooth projective space. Complete
presentations prepare an algebraic model stable under rescaling.
\item \textbf{Uniform division and $\bar\partial$ lifting.}
A weighted Gr\"obner-basis argument gives uniform local division.
Sampled Hermitian metrics control differentiated generators and
curvature in compatible norms. A finite correction through the
syzygies, using $\bar\partial$ on the presenting vector bundles,
produces the global holomorphic lift.
\item \textbf{Return to the shift.}
Differentiation and projection turn the lift into the scalar
coordinate-defect estimate. Auxiliary variables transfer that bound
to the original vector-valued module. Its Hilbert function then
determines the Schatten exponent.
\end{enumerate}
The geometric core below is scalar, as in the full manuscript.
Step~3 recovers finite coefficient multiplicity by an explicit
identity before singular values are counted.

\section{Step 1: projective realization of the graded problem}
\label{sec:step-one}

We first work with an arbitrary scalar homogeneous ideal
$I\subset R^{(d)}=\C[x_1,\ldots,x_d]$. The number $d$ may later be
$m+r$ in the auxiliary-variable construction. Put $N=d-1$,
$W=\PP^N$, and $L=\cO_W(1)$ with its Fubini--Study metric.
The saturation $I^{\mathrm{sat}}/I$ is finite dimensional, and
\[
 R^{(d)}_n=H^0(W,L^n),\qquad
 I_n^{\mathrm{sat}}=H^0(W,\cI\otimes L^n),\qquad
 \cI=\widetilde I\subset\cO_W.
\]
Thus $I_n=H^0(W,\cI\otimes L^n)$ for all sufficiently large $n$.
For normalized Fubini--Study volume the norm correspondence is exact:
\begin{equation}\label{eq:norm}
 \|F_n\|_{\DA}^2=c_n\|F_n\|_{\FS}^2,
 \qquad c_n=\binom{n+d-1}{d-1},\qquad
 \frac{c_n}{c_{n+1}}=\frac{n+1}{n+d}.
\end{equation}
Indeed, monomials are orthogonal and
$\|x^\alpha\|_{\FS}^2=(d-1)!\alpha!/(n+d-1)!$ for $|\alpha|=n$.
Consequently degreewise orthogonal projections agree in the two
realizations. The finitely many degrees changed by saturation do not
affect Schatten membership. This transfer uses homogeneity to
preserve both the relations and the Hilbert-space geometry.

The cases $d=1$ and the zero or unit ideal sheaf are elementary.
Assume henceforth that $N\ge1$ and $\cI$ is nonzero and proper.

\subsection{Theorem B: preserving the relations under rescaling}
The local analysis must remain uniform when the center moves and the
observation scale tends to zero. Normalizing the generators prevents
their coefficients from vanishing, but by itself gives no control of
the relations among their polynomial multiples. The construction
therefore chooses complete spaces of sections at every syzygy stage. Choose degrees $0\le D_0\le\cdots\le D_N$ and a complex
\begin{equation}\label{eq:complex}
 E_N\xrightarrow{d_N}\cdots\xrightarrow{d_1}E_0
 \xrightarrow{d_0=\Phi}\cO_W,
 \qquad E_q=\mathcal V_q\otimes L^{-D_q},
\end{equation}
by setting $\cK_0=\cI$ and, recursively,
\[
 \mathcal V_q=H^0(W,\cK_q(D_q)),\qquad
 \im d_q=\cK_q,\qquad \cK_{q+1}=\ker d_q.
\]
\begin{remark}[The resolution antecedent]
The resolution viewpoint has a direct antecedent in the exact
Hilbert-module complexes of Douglas--Jabbari--Tang--Yu
\cite{DouglasJabbariTangYu}. Here \eqref{eq:complex}, with its last
map factored through $\cI$, is a partial locally free resolution of a
coherent ideal sheaf: it is exact at $E_0,\ldots,E_{N-1}$, and
injectivity of $d_N$ is not asserted. Its successive syzygies correct
holomorphic overlap errors. Sheaf exactness alone does not ensure
exactness after Hilbert-space completion; the uniform division and
metric estimates below supply the analytic control needed for the
lifting theorem.
\end{remark}

Here $D_0$ is regular uniformly over the Hilbert scheme with the
prescribed quotient polynomial; each subsequent $D_q$ is regular for
every kernel in $\Quot_W(E_{q-1},P_{\cK_{q-1}})$, the proper Quot
family with the indicated quotient Hilbert polynomial.
Such degrees exist by uniform Castelnuovo--Mumford regularity
\cite{Stacks}. Using the complete section space, in a degree chosen
for the entire family, is essential.

Write $B_s=\{z\in\C^N:|z|<s\}$. On a fixed buffered affine chart define
\begin{equation}\label{eq:gram}
 G_{-1}=1,\qquad
 G_q(p,\eps)=\int_{B_1}d_q(p+\eps z)^*G_{q-1}(p,\eps)
                    d_q(p+\eps z)\,d\mu(z),
\end{equation}
where $\mu$ is probability volume. At every nonzero stage $G_q$ is
positive definite. Set
\begin{equation}\label{eq:normalized}
 A_q^{p,\eps}(z)=G_{q-1}(p,\eps)^{1/2}
 d_q(p+\eps z)G_q(p,\eps)^{-1/2}.
\end{equation}
The normalization gives
\[
\int_{B_1}(A_q^{p,\eps})^*A_q^{p,\eps}\,d\mu=\Id.
\]
Thus the columns form an orthonormal frame in a fixed
finite-dimensional polynomial space, proving coefficient compactness.
The next theorem additionally asserts preservation of the images
and kernels, which requires the complete-section construction.


\begin{theorem}[Complete presentations under normalized limits]
\label{thm:compact}
The presentation \eqref{eq:complex} can be chosen so that the tuples
$\bigl(A_0^{p,\eps},\ldots,A_N^{p,\eps}\bigr)$, with centers in
fixed relatively compact subcharts of a finite buffered affine atlas,
and sufficiently small $\eps>0$,
have compact closure in coefficient topology. After homogenization
to the prescribed degrees $D_q-D_{q-1}$, every limiting tuple is
exact as a sheaf complex at $E_0,\ldots,E_{N-1}$.

Moreover, the homogeneous polynomial image $J_q$ of each normalized
or limiting map
\[
 \C[Z_0,\ldots,Z_N](-D_q)^{\dim\mathcal V_q}
 \longrightarrow
 \C[Z_0,\ldots,Z_N](-D_{q-1})^{\dim\mathcal V_{q-1}}
\]
has the same Hilbert function in every degree. With $\mathcal V_{-1}=\C$,
$E_{-1}=\cO_W$, $D_{-1}=0$, and $P_q$ the Hilbert polynomial of
$\cK_q$, it is
\begin{equation}\label{eq:hilbert}
 \dim(J_q)_\ell=
 \begin{cases}0,&\ell<D_q,\\ P_q(\ell),&\ell\ge D_q.\end{cases}
\end{equation}
\end{theorem}

The normalization gives $\int_{B_1}(A_q^{p,\eps})^*A_q^{p,\eps}\,d\mu=\Id$;
the columns therefore lie in a compact space of orthonormal frames of
fixed-degree polynomials. Preserving their polynomial multiples is
the deeper point. Affine dilations extend to projective automorphisms, and each
constant coefficient change acts on an equally twisted bundle.
The transformed kernels are thus points of fixed proper Quot
schemes. The universal kernels are flat over these parameter spaces;
uniform regularity and cohomology and base change identify their
degree-$D_q$ sections with a vector bundle \cite{Stacks,Conrad}.
Its inclusion in the fixed space
$H^0(W,E_{q-1}(D_q))$ defines a Grassmannian map with closed image.
Every limiting column space remains the \emph{complete} section
space of a kernel in that family. It generates that kernel, and
regularity identifies all its higher multiples with the higher
section spaces.


Finally, the complex identities pass to the limit. Inductively, the
limiting next image is contained in the preceding kernel, and the
two sheaves have the same Hilbert polynomial. Their quotient has
zero Hilbert polynomial and hence is the zero coherent sheaf.
This proves exactness and \eqref{eq:hilbert}; below $D_q$ the
polynomial image is zero by its generating degree. Theorem B applies
classical properness, regularity, and base change to the linked
normalizations of an entire presentation. Its role is to give
Theorem C a compact family in which image dimensions do not jump.

\section{Step 2: uniform division and global holomorphic lifting}
\label{sec:step-two}

The analytic task is to solve $\Phi a=F$ with a holomorphic
coefficient section $a$ whose norm is controlled uniformly in the
large degree. The operators $\dbar$ act on forms with values in
$E_q\otimes L^h$. The ideal sheaf records the relations; its locally
free presentation supplies the bundles on which the estimates are
made. The three ingredients are uniform division, compatible metrics,
and a finite global correction.

\subsection{Theorem C: division  with constants uniform over the family}
For a homogeneous matrix $A$, write $A_{\mathrm{aff}}(z)=A(1,z)$.
Membership in its analytic image sheaf means local holomorphic
divisibility near every point under consideration.

\begin{theorem}[Uniform holomorphic division]\label{thm:division}
Let $\mathcal A$ be a compact coefficient family of homogeneous
matrices
\[
 A:\C[Z_0,\ldots,Z_N](-e)^a\longrightarrow
   \C[Z_0,\ldots,Z_N]^b
\]
of fixed size and degree $e\ge0$. Suppose that $\dim(\im A)_\ell$ is
independent of $A\in\mathcal A$ for every $\ell$.
For every $s>0$ there are $R_s>s$ and $C_s<\infty$ such that, for every
$A\in\mathcal A$, every square-integrable holomorphic $\C^b$-valued
$F$ belonging to its analytic image sheaf throughout $B_{R_s}$
admits a holomorphic $\C^a$-valued lift $u$, linear in $F$, on $B_s$ with
\begin{equation}\label{eq:division}
 A_{\mathrm{aff}}u=F|_{B_s},\qquad
 \|u\|_{L^2(B_s)}\le C_s\|F\|_{L^2(B_{R_s})}.
\end{equation}
\end{theorem}

Galligo's parameter-dependent vector division theorem controls
quotients and remainders in weighted coefficient norms; his
flatness criterion identifies when the remainder vanishes on the
module \cite[Theorems~1.2.7--1.2.8]{Galligo}. His
Theorem~1.4.2 gives parameter-linear splittings. Hauser--M\"uller
construct compatible continuous linear splittings on weighted
Banach scales, with explicit bounds of the form
$\|\sigma_s\|\le C s^{-\beta}$
\cite[Theorems~5.1--5.2]{HauserMuller}.
Theorem C adapts this quantitative division to compact families of
homogeneous matrices and states the conclusion in two-ball $L^2$
norms. Analytic dependence of the lift on a parameter is not part of
its assertion. The weighted division and perturbation mechanism
already occur in these classical works.


Use one homogeneous Gr\"obner basis at a reference matrix.
Choose a weighted coefficient $\ell^1$ norm so that each leading
monomial dominates its tail. The monomial decomposition is an
isometry on one Banach space containing all degrees. Nearby matrices
give a perturbation of norm at most $1/2$, hence a degree-preserving
Neumann inverse bounded by $2$. Equality of the image dimensions
forces the division remainder to vanish on the image in each degree.
The homogenizing variable is assigned weight one, so dehomogenization
incurs no loss depending on the representation degree.

For holomorphic data, Cartan's theorem~B supplies an initial,
unestimated preimage \cite{GrauertRemmert}. Approximating that preimage
produces polynomial \emph{image} approximants, to which the bounded
division is applied afresh. Cauchy estimates and a finite compactness
cover give \eqref{eq:division}. The theorem provides a sufficiently
large input ball for each output ball; its hypothesis concerns the
entire input ball. Theorem~\ref{thm:compact} verifies its assumptions
for all stages of the presentation, after removing the common
target grading shift.


The full Hilbert-function hypothesis is convenient because it gives
division in every homogeneous degree. For the affine conclusion
alone, a common eventual Hilbert polynomial suffices: a polynomial
image datum may be homogenized beyond the relevant stabilization
degrees, and the unit weight of the homogenizing variable makes the
estimate independent of that choice. The stronger formulation above
is exactly the one supplied by Theorem B.


The equation in Theorem C lifts an ambient holomorphic datum through
a presentation. In a restriction--extension theorem, the datum is
instead prescribed on a subvariety. The latter is the setting of
Beatrous' estimates \cite{Beatrous} and their extension in Douglas--Tang--Yu
\cite[Theorem~4.3]{DouglasTangYu}; it does not by itself provide
the family of division operators asserted here.

\subsection{Theorem D: a metric adapted to division and differentiation}
The coefficient metric measures the size of a lift. It is auxiliary:
the scalar datum and the final Drury--Arveson norm retain their
original metrics. This freedom is necessary even for a smooth
principal ideal. On $\PP^1$, let
$\Phi=Z_0:\cO(-1)\to\cO$ and $F_n=Z_0Z_1^n$.
The unique lift is $a_n=Z_1^n$, and normalized Fubini--Study norms give
\[
 \|a_n\|^2=\frac{1}{n+1},\qquad
 \|F_n\|^2=\frac{1}{(n+1)(n+2)}.
\]
The lifting norm therefore grows as $\sqrt{n+2}$ in the standard
coefficient metric, and comparison with any fixed smooth positive
metric cannot remove this growth. The illustrative choice
\[
|v|_{H_\eps}^2=|\Phi v|^2+\eps^2|v|_{\mathrm{std}}^2
\]
instead gives
$\|a_n\|_{H_\eps\otimes L^{n+1}}^2/\|F_n\|_{L^{n+1}}^2
=1+\eps^2(n+2)$, which stays bounded when $\eps$ is a fixed
multiple of $n^{-1/2}$. For a general presentation, one needs this
adaptation together with bounds for all differentials, curvature,
and differentiated generators.


Sampled norms are scaled jet norms. At the first stage, orthogonality of monomials on $B_1$ gives the
exact identity
\begin{equation}\label{eq:jets}
 G_0(p,\eps)[v]
=\sum_{|\alpha|\le D_0}
 \frac{N!\,\eps^{2|\alpha|}}{\alpha!(N+|\alpha|)!}
 |(\partial^\alpha\Phi)(p)v|^2.
\end{equation}
The metric records all the finitely many jets, with their natural
powers of the sampling radius. A coefficient direction remains
visible even where its value vanishes, and one derivative costs a
factor $\eps^{-1}$.

\subsection{Sampling the presentation}
%DIFAUXCMD
\addtocounter{subsection}{-1}%DIFAUXCMD
\DIFdelend The recursion \eqref{eq:gram} can be expanded as
\begin{equation}\label{eq:independent}
 G_q(p,\eps)[v]=\int_{(B_1)^{q+1}}
 \bigl|d_0(p+\eps z_0)\cdots d_q(p+\eps z_q)v\bigr|^2\,d\mu^{q+1}.
\end{equation}
The sample points are independent: evaluation of all factors at one
point would give zero by the complex identities. Positivity follows
inductively. If the integral defining $G_q[v]$ vanishes, positivity
of $G_{q-1}$ forces $d_q(\cdot)v$ to be the zero polynomial. Independence of the columns over constants then gives $v=0$.

\begin{theorem}[Compatible metrics with controlled curvature]
\label{thm:metrics}
For the fixed complete presentation \eqref{eq:complex}, there are smooth
positive Hermitian metrics $H_{q,\eps}$, $0<\eps\le\eps_0$, with
$H_{-1,\eps}=1$, such that every $d_q$ is contractive and
\begin{equation}\label{eq:metric-estimates}
 \|\Theta(H_{q,\eps})\|_{H_{q,\eps},\omega_{\FS}}
       \le C\eps^{-2},\qquad
 |(\partial_{x_j}\Phi)v|_{L^{-1}}
       \le C\eps^{-1}|v|_{H_{0,\eps}}.
\end{equation}
The local metrics are uniformly comparable to the sampled forms
$G_q$ on their actual bundle fibers. The derivative is taken in an
ambient homogeneous coordinate and defines a map $E_0\to L^{-1}$;
the target carries its original Fubini--Study metric.
\end{theorem}

View the product polynomial in \eqref{eq:independent} as a map
$F_q(p)$ into a fixed finite-dimensional Hilbert space. Then
$G_q=F_q^*F_q$ and
\[
 \partial_{p_a}F_q
 =\eps^{-1}\sum_{j=0}^q\partial_{z_{j,a}}F_q.
\]
Differentiation is bounded on its fixed finite-dimensional polynomial
space. This identity controls the first and second derivatives of
$G_q$ \emph{relative to $G_q$} by $C\eps^{-1}$ and $C\eps^{-2}$,
even as eigenvalues in a fixed coefficient frame approach zero.
The comparison concerns $G_q^{-1/2}(\partial G_q)G_q^{-1/2}$ and
the corresponding second derivatives; no uniform Euclidean lower
bound for $G_q$ is needed. On overlapping charts the forms are
compared on the same bundle  fibers and patched to metrics
$\widehat G_q$. 

Define the graph metrics recursively by
\[
H_q=\widehat G_q+d_q^*H_{q-1}d_q
     =\widehat G_q+d_q^*\widehat G_{q-1}d_q,
\]
where $\widehat G_{-1}=1$ and the second equality follows from
$d_{q-1}d_q=0$ for $q\ge1$ (and directly from the definition for
$q=0$). These metrics make all differentials contractive.
The Chern curvature formula
$\Theta(H)=\dbar(H^{-1}\partial H)$ converts the relative first
and second derivative bounds into $O(\eps^{-2})$
\cite[Section~3]{Demailly}. Derivative evaluation in
\eqref{eq:jets}, followed by  Euler's homogeneous identity, gives the second estimate in \eqref{eq:metric-estimates}.

Hermitian splittings, curvature estimates, and $L^2$ lifting already
play a central role in Skoda's bundle-morphism theorem
\cite{Skoda1978}. The assertion needed here is the simultaneous
control of division, differentiation, and curvature by these sampled
metrics. Positive definiteness of a metric does not imply positive
curvature; its possible negative curvature will be dominated by
the growing positive line-bundle twist. The construction of D uses
the fixed polynomial presentation and constant independence of its
columns. The compactness and Hilbert-function conclusions of B
supply the uniform division input. Exactness of the actual
presentation places holomorphic cycles in the required image sheaves.

\subsection{Theorem E: bounded projective lifting}
All $L^2$ norms below use normalized Fubini--Study volume and the displayed
coefficient and line metrics.

\begin{theorem}[Bounded projective lifting]\label{thm:lift}
There are fixed $R_*>0$, $n_0$, and $C$ such that, for
$n\ge n_0$, $\eps_n=R_*/\sqrt n$, and $h\in\{n,n+1\}$, there is
a linear map
\[
 T_{n,h}:H^0(W,\cI\otimes L^h)
      \longrightarrow H^0(W,E_0\otimes L^h)
\]
satisfying
\begin{equation}\label{eq:lift}
 \Phi T_{n,h}F=F,\qquad
 \|T_{n,h}F\|_{H_{0,\eps_n}\otimes L^h}
       \le C\|F\|_{L^h}.
\end{equation}
The datum retains its original scalar norm, and the same coefficient
metric serves both adjacent line powers.
\end{theorem}

\emph{Projective lifting} means finding a global holomorphic
coefficient section above the datum under $\Phi$.
Serre vanishing gives surjectivity in sufficiently high degree;
the additional assertion in E is the bound uniform in $n$, in the
metrics of D and with the datum's original scalar norm.


On balls of radius $O(\eps)$, a holomorphic line
frame removes the constant and linear terms of the Fubini--Study
potential. The remaining weight variation is bounded by
$\exp(C h\eps^2)$. Theorems~\ref{thm:division} and
\ref{thm:metrics} therefore give local holomorphic division with
uniform constants whenever $h\eps^2$ is bounded.

%DIFDELCMD < %%%
\DIFdelend On $E_q\otimes L^h$, the Bochner--Kodaira--Nakano estimate,
applied in top holomorphic degree after the canonical $K_W^{-1}$
twist, gives a solution operator for closed positive-degree forms with
\begin{equation}\label{eq:dbar}
 \dbar u=\zeta,\qquad
 \|u\|\le(c_Lh-C_0-C_1\eps^{-2})^{-1/2}\|\zeta\|,
\end{equation}
when the quantity in parentheses is positive \cite{Demailly}.
Choose $R_*$ large enough that $C_1/R_*^2<c_L/2$, and then choose
$n_0$ large. At $\eps=\eps_n$ and $h=n,n+1$, the solution operator
has norm $O(\eps_n)$. Thus the two requirements are compatible:
the local line weight has variation $\exp(O(h\eps_n^2))$,
while the positive curvature of $L^h$ dominates the
$O(\eps_n^{-2})$ coefficient-curvature loss. The constant $R_*$ is
fixed first; it need not be small.


Take a bounded-overlap cover at this scale with buffered balls.
Local lifts and their holomorphic alternating differences are divided
successively through $d_1,\ldots,d_N$. In \v Cech notation, these
cochains satisfy
\[
\Phi c_0=F,\qquad d_qc_q=\delta c_{q-1}.
\]
The right side is a holomorphic cycle because $\delta^2=0$ and
the preceding lifting identity holds. Exactness places it in the
next image sheaf, where C applies. The buffers keep every such
datum inside the entire input ball required for division. Only then introduce a partition of unity and set
\[
 \eta_q=\sum_{i_0,\ldots,i_q}
 \chi_{i_0}\dbar\chi_{i_1}\wedge\cdots\wedge\dbar\chi_{i_q}
 c_{q,i_0\ldots i_q}.
\]
The ordered alternating-cochain convention gives
\[
\Phi\eta_0=F,\qquad
 d_q\eta_q=\dbar\eta_{q-1},\qquad
 \dbar\eta_N=0,\qquad
 \|\eta_q\|\le C_q\eps_n^{-q}\|F\|.
\]
\DIFaddend Bounded overlap controls the sum of local squared norms, independently
of the number of balls. The factor $\eps_n^{-q}$ comes from the
$q$ differentiated cutoffs. All resulting forms take values in the
smooth bundles $E_q\otimes L^h$; no kernel sheaf is assumed to be a
vector bundle.

Now set $\zeta_N=\eta_N$ and descend by
\[
 \dbar u_q=\zeta_q,\qquad
 \zeta_{q-1}=\eta_{q-1}-d_qu_q,
 \qquad q=N,\ldots,1.
\]
The top form is closed for dimensional reasons. At each lower stage,
\[
\dbar\zeta_{q-1}
 =d_q(\eta_q-\zeta_q).
\]
This is zero when $q=N$ because $\zeta_N=\eta_N$; when $q<N$,
it equals $d_qd_{q+1}u_{q+1}=0$. Each solution
recovers one factor $\eps_n$ lost to a differentiated cutoff, and
contractivity of $d_q$ preserves the estimate. The finite recursion
yields a holomorphic $a=\zeta_0$ with $\Phi a=F$ and
$\|a\|\le C\|F\|$. The presentation, curvature constants and
division radii are fixed before $R_*$ is chosen. Local constants may
grow with $R_*$; a finite descent requires only that they be fixed
as $n$ varies. More explicitly, if the $\dbar$ solver has norm at most
$b\eps_n$ and $\|\eta_q\|\le C_q\eps_n^{-q}\|F\|$, induction gives
\[
\|a\|\le \left(\sum_{q=0}^N b^qC_q\right)\|F\|.
\]
There is no convergence of an infinite correction scheme to prove.
The same metrics serve $h=n,n+1$; the two linear lifting operators
may be chosen separately.

H\"ormander's work on generators and Kelleher--Taylor's weighted
ideal theory use the interaction between holomorphic division and
$\dbar$ \cite{HormanderGenerators,KelleherTaylor}.
Kelleher--Taylor explicitly organize the correction with a
Koszul--Dolbeault complex \cite[Section~II]{KelleherTaylor}.
Skoda's $L^2$ division theorem obtains quantitative holomorphic
coefficients under generator-dependent integrability conditions
\cite{Skoda1972}; its proof uses a direct estimate rather than
that Koszul construction. His later bundle-morphism theorem starts
from a smooth Hermitian splitting and corrects its $\dbar$ error
\cite[Section~2]{Skoda1978}.


Andersson uses descending $\dbar$ correction in polynomial division
\cite[Theorem~2.3 and Example~2.4]{Andersson2006}.
Andersson--Wulcan then work with general locally free complexes and
residue currents. In their projective application, a graded free
resolution is sheafified, twisted by $\cO(r)$, and used in successive
global $\dbar$ equations \cite[Section~6, Proposition~6.1]{AnderssonWulcan}.
This is a direct structural antecedent of E. The present construction
starts with local holomorphic divisions and forms obtained from
\v Cech cochains. The quantitative task is to keep their correction
bounded in the degree-dependent metrics, while retaining the
original norm of every datum in the image sheaf.


Varolin's twisted division estimates and Ji's division theorems for
generically exact bundle sequences further develop the role of
curvature and $L^2$ estimates in division
\cite{Varolin,Ji}. Their hypotheses and weighted norms must be
compared with the desired estimate: E concerns every projective ideal
section in its original scalar norm. Its specific claim is
simultaneous uniformity in the line power, local scale, and
normalized matrix families.


\section{Step 3: return to the Drury--Arveson shift}
\label{sec:step-three}

\subsection{Differentiation and projection}
For the scalar ideal from Step~1, write
$Q_n^I=R_n^{(d)}\ominus I_n$, $X_j=M_{x_j}$, and
$B_{j,n}^I=P_{I_{n+1}}X_j|_{Q_n^I}$. For $n\ge n_0$ with
both $n$ and $n+1$ beyond the saturation discrepancy, take
$F\in I_{n+1}$ and let
$a=T_{n,n+1}F$. Differentiating the holomorphic lift $F=\Phi a$ gives
\[
 P_{Q_n^I}\partial_{x_j}F
 =P_{Q_n^I}\bigl((\partial_{x_j}\Phi)a\bigr),
\]
because $\Phi\partial_{x_j}a\in I_n$. Thus projection eliminates
derivatives of the coefficient section. On degree $n+1$,
$X_j^*=(n+1)^{-1}\partial_{x_j}$. Theorems~\ref{thm:metrics}
and~\ref{thm:lift}, followed by the exact norm conversion
\eqref{eq:norm}, give
\begin{align}
\|(B_{j,n}^I)^*F\|_{\DA}
 &\le\frac{C}{\eps_n(n+1)}
       \sqrt{\frac{c_n}{c_{n+1}}}\,\|F\|_{\DA}\notag\\
 &=\frac{C}{\eps_n\sqrt{(n+1)(n+d)}}\|F\|_{\DA}
   =O(n^{-1/2})\|F\|_{\DA}.\label{eq:return}
\end{align}
Summing the squared bounds and absorbing the finitely many omitted
blocks proves the scalar estimate in every ambient dimension:
\begin{equation}\label{eq:scalar-defect}
 A_n^I=\sum_{j=1}^d(B_{j,n}^I)^*B_{j,n}^I
 \le\frac{C_I}{n+1}\Id_{Q_n^I}.
\end{equation}

\subsection{Finite coefficient multiplicity}
Return to $M\subset R\otimes\C^r$ from Theorem~\ref{thm:main}.
Choose an orthonormal coefficient basis and homogeneous generators
$f^{(1)},\ldots,f^{(s)}$ of $M$. Adjoin $r$ variables and set
\begin{equation}\label{eq:encoding}
 \lambda(f)=\sum_{b=1}^r f_b(z)y_b,\qquad
 J_M=(\lambda(f^{(1)}),\ldots,\lambda(f^{(s)}))+(y_1,\ldots,y_r)^2.
\end{equation}
This is an ideal in $R'=R[y_1,\ldots,y_r]$, homogeneous in both
groups of variables separately. Its part of $y$-degree one is
$\lambda(M)$; all higher $y$-degrees lie in $J_M$. The encoding and
its degreewise normalization follow Shalit's auxiliary-variable
method \cite[Section~5, equation~(5.1)]{ShalitDivision}. The
square-zero relations make the complete positive defect explicit.

Let $\mathcal E_n$ be the sector of the scalar quotient by
$\overline{J_M}$ having $z$-degree $n$ and $y$-degree one. Its total
degree is $n+1$. The map
\begin{equation}\label{eq:unitary}
 U_n:Q_n\longrightarrow\mathcal E_n,\qquad
 U_nq=\sqrt{n+1}\sum_{b=1}^r q_b(z)y_b
\end{equation}
is unitary: the squared monomial norms are respectively
$\alpha!/n!$ and $\alpha!/(n+1)!$, and the map identifies the
corresponding polynomial submodules.

Let $A_\ell^{J_M}$ be the scalar defect summed over all $m+r$
coordinates. The exact transfer identity is
\begin{equation}\label{eq:vector-transfer}
 U_n^*A_{n+1}^{J_M}U_n
 =\frac{n+1}{n+2}A_n+\frac{r+1}{n+2}\Id_{Q_n}.
\end{equation}
Indeed, multiplication by $z_j$ changes the unitary normalization by
$\sqrt{(n+1)/(n+2)}$. Every product $y_bU_nq$ belongs to $J_M$,
and the sum of its squared forward weights over $b$ is
$(r+1)/(n+2)$. The bigrading makes the auxiliary-degree sectors
orthogonal. Applying \eqref{eq:scalar-defect} in total degree $n+1$
and dropping the positive scalar term gives
\[
 \frac{n+1}{n+2}A_n\le\frac{C_{J_M}}{n+2}\Id,
 \qquad A_n\le\frac{C_{J_M}}{n+1}\Id.
\]
This proves the defect estimate in the original vector-valued norm.
Equations~\eqref{eq:compression}--\eqref{eq:ambient} then give the
quotient commutator estimate in Theorem~\ref{thm:main}.

\subsection{Counting in the original quotient}
For infinite-dimensional $T=(R\otimes\C^r)/M$ and
$\delta=\dim\Supp_R T$, the Hilbert polynomial gives
$\dim Q_n=\dim T_n=O((n+1)^{\delta-1})$. A block of norm at most
$C/(n+1)$ contributes no singular values greater than $t$ when
$n+1>C/t$. Consequently
\[
 \#\{k:s_k([S_i^*,S_j])>t\}
 \le\sum_{n+1\le C/t}\dim Q_n=O(t^{-\delta}).
\]
This is \eqref{eq:weak}. Counting the column defect, whose
degree-$n$ norm is $O((n+1)^{-1/2})$, gives the exponent $2\delta$.
Each $[P_{\mathcal M},Z_j]$ is the corresponding off-diagonal block
extended by zero. For the restriction tuple,
\[
 [V_i^*,V_j]|_{M_n}
 =P_{M_n}[Z_i^*,Z_j]|_{M_n}+B_{j,n-1}B_{i,n-1}^*,
\]
where the last term is zero at $n=0$. These blocks are
$O((n+1)^{-1})$ and
$\dim M_n\le r\binom{n+m-1}{m-1}$, giving the weak exponent $m$.
If $T$ is finite dimensional, the quotient and off-diagonal terms
have finite rank, while the restriction estimate remains valid.

Since $\sqrt{J_M}=(y_1,\ldots,y_r)$, the auxiliary quotient has
dimension $m$, possibly larger than $\delta$. Counting there loses the sharper exponent;
 replacing $J_M$ by its radical deletes the
sector encoding $T$.



\begin{thebibliography}{99}
% Alphabetized by author surname; anonymous companion manuscript by title.
% Compared by citation key; reordering alone is not marked. Deleted entries remain in red.

\bibitem{Andersson2006}
M.~Andersson, The membership problem for polynomial ideals in
terms of residue currents, \emph{Ann. Inst. Fourier (Grenoble)}
\textbf{56} (2006), no.~1, 101--119.
\href{https://doi.org/10.5802/aif.2174}{doi:10.5802/aif.2174}.

\bibitem{AnderssonWulcan}
M.~Andersson and E.~Wulcan, Residue currents with prescribed
annihilator ideals, \emph{Ann. Sci. \'Ecole Norm. Sup. (4)}
\textbf{40} (2007), no.~6, 985--1007.
\href{https://doi.org/10.1016/j.ansens.2007.11.001}{doi:10.1016/j.ansens.2007.11.001}.

\bibitem{Arveson1998}
W.~Arveson, Subalgebras of $C^*$-algebras III: multivariable operator
theory, \emph{Acta Math.} \textbf{181} (1998), 159--228.
\url{https://doi.org/10.1007/BF02392585}.


\bibitem{Arveson2005}
W.~Arveson, $p$-summable commutators in dimension $d$,
\emph{J. Operator Theory} \textbf{54} (2005), 101--117.
\url{https://www.jstor.org/stable/24715673}.


\bibitem{Beatrous}
F.~Beatrous, Jr., $L^p$-estimates for extensions of holomorphic
functions, \emph{Michigan Math. J.} \textbf{32} (1985), no.~3,
361--380.
\href{https://doi.org/10.1307/mmj/1029003244}{doi:10.1307/mmj/1029003244}.

\bibitem{Conrad}
B.~Conrad, \emph{Applications of base change for coherent cohomology},
Math 248B notes, Theorem~1.1 and Corollary~1.2.
\url{https://math.stanford.edu/~conrad/248BPage/handouts/cohom.pdf}.

\bibitem{Demailly}
J.-P.~Demailly, \emph{Applications of the theory of $L^2$ estimates
and positive currents in algebraic geometry}, lecture notes, 2007,
Sections~3, 5--6, and~12. 
\url{https://www-fourier.univ-grenoble-alpes.fr/~demailly/manuscripts/eem2007.pdf}.

\bibitem{Ding26}
L.~Ding, {Essential normality of holomorphic Sobolev quotient submodules}, \emph{Trans. Amer. Math. Soc.} \textbf{379} (2026), 7525–7555. \url{https://doi.org/10.1090/tran/9717}.


\bibitem{DingFangTang}

L.~Ding, X.~Fang, and X.~Tang,
\emph{Projective division and essential normality of graded Hilbert modules},
companion research manuscript, 20 November 2026, 39 pp.;
Theorem~1.1 and Sections~3--8. Working draft for independent
mathematical assessment.


\bibitem{DouglasIndex}
R.~G.~Douglas, A new kind of index theorem, in
\emph{Analysis, Geometry and Topology of Elliptic Operators},
World Scientific, 2006, 369--382.
\url{https://doi.org/10.1142/9789812773609_0014}.

\bibitem{DouglasJabbariTangYu}
R.~G.~Douglas, M.~Jabbari, X.~Tang, and G.~Yu,
A new index theorem for monomial ideals by resolutions,
\emph{J. Funct. Anal.} \textbf{275} (2018), no.~3, 735--760. \href{https://doi.org/10.1016/j.jfa.2018.04.003}{doi:10.1016/j.jfa.2018.04.003}.
\url{https://arxiv.org/abs/1708.05796}.

\bibitem{DouglasTangYu}
R.~G.~Douglas, X.~Tang, and G.~Yu, An analytic Grothendieck Riemann
Roch theorem, \emph{Adv. Math.} \textbf{294} (2016), 307--331.
\href{https://doi.org/10.1016/j.aim.2016.02.031}{doi:10.1016/j.aim.2016.02.031}.
\url{https://arxiv.org/abs/1404.4396}.

\bibitem{EnglisEschmeier}
M.~Engli\v{s} and J.~Eschmeier, Geometric Arveson--Douglas conjecture,
\emph{Adv. Math.} \textbf{274} (2015), 606--630; corrigendum,
\emph{Adv. Math.} \textbf{278} (2015), 254.
\href{https://doi.org/10.1016/j.aim.2014.11.026}{doi:10.1016/j.aim.2014.11.026}.


\bibitem{EschmeierCriterion}
J.~Eschmeier, Essential normality of homogeneous submodules,
\emph{Integral Equations Operator Theory} \textbf{69} (2011), 171--182.
\href{https://doi.org/10.1007/s00020-010-1809-4}{doi:10.1007/s00020-010-1809-4}.

\bibitem{Galligo}
A.~Galligo, Th\'eor\`eme de division et stabilit\'e en g\'eom\'etrie
analytique locale, \emph{Ann. Inst. Fourier (Grenoble)}
\textbf{29} (1979), no.~2, 107--184.
\href{https://doi.org/10.5802/aif.745}{doi:10.5802/aif.745}.

\bibitem{GrauertRemmert}
H.~Grauert and R.~Remmert, \emph{Theory of Stein Spaces},
Grundlehren der mathematischen Wissenschaften, vol.~236,
Springer, 1979.

\bibitem{GuoWang}
K.~Guo and K.~Wang, Essentially normal Hilbert modules and
$K$-homology, \emph{Math. Ann.} \textbf{340} (2008), 907--934.
\href{https://doi.org/10.1007/s00208-007-0175-2}{doi:10.1007/s00208-007-0175-2}.


\bibitem{HauserMuller}
H.~Hauser and G.~M\"uller, A rank theorem for analytic maps between
power series spaces, \emph{Publ. Math. Inst. Hautes \'Etudes Sci.}
\textbf{80} (1994), 95--115.
\href{https://doi.org/10.1007/BF02698897}{doi:10.1007/BF02698897}.

%\bibitem{HeltonHowe75}
%{\color{red}
%J.~Helton, R.~Howe, {Traces of commutators of integral operators,} \emph{Acta Math.} \textbf{135} (1975), 271-305. \url{https://doi.org/10.1007/BF02392022}.
%}

\bibitem{HormanderGenerators}
L.~H\"ormander, Generators for some rings of analytic functions,
\emph{Bull. Amer. Math. Soc.} \textbf{73} (1967), 943--949.
\href{https://doi.org/10.1090/S0002-9904-1967-11860-3}{doi:10.1090/S0002-9904-1967-11860-3}.

\bibitem{Ji}
Q.~Ji, Division theorems for exact sequences,
\emph{Math. Ann.} \textbf{355} (2013), no.~3, 947--984.
\href{https://doi.org/10.1007/s00208-012-0807-z}{doi:10.1007/s00208-012-0807-z}.

\bibitem{KelleherTaylor}
J.~J.~Kelleher and B.~A.~Taylor, Finitely generated ideals in rings
of analytic functions, \emph{Math. Ann.} \textbf{193} (1971),
225--237.
\href{https://doi.org/10.1007/BF02052394}{doi:10.1007/BF02052394}.

\bibitem{Companion}
\emph{Projective division and essential normality of graded Hilbert modules},
companion research manuscript, 20 November 2026, 39 pp.;
Theorem~1.1 and Sections~3--8. Working draft for independent
mathematical assessment.

\bibitem{ShalitDivision}
O.~M.~Shalit, Stable polynomial division and essential normality of
graded Hilbert modules, \emph{J. Lond. Math. Soc. (2)}
\textbf{83} (2011), 273--289.
\url{https://doi.wiley.com/10.1112/jlms/jdq054}.


\bibitem{Skoda1972}
H.~Skoda, Application des techniques $L^2$ \`a la th\'eorie des
id\'eaux d'une alg\`ebre de fonctions holomorphes avec poids,
\emph{Ann. Sci. \'Ecole Norm. Sup. (4)} \textbf{5} (1972), no.~4,
545--579.
\href{https://doi.org/10.24033/asens.1237}{doi:10.24033/asens.1237}.

\bibitem{Skoda1978}
H.~Skoda, Morphismes surjectifs de fibr\'es vectoriels semi-positifs,
\emph{Ann. Sci. \'Ecole Norm. Sup. (4)} \textbf{11} (1978), no.~4,
577--611.
\href{https://doi.org/10.24033/asens.1357}{doi:10.24033/asens.1357}.

\bibitem{Stacks}
The Stacks Project Authors, \emph{The Stacks Project},
Castelnuovo--Mumford regularity, Tags 08A6--08A8 and 08AG;
properness of Quot, Tag 0DPA.
\url{https://stacks.math.columbia.edu/tag/08AG};
\url{https://stacks.math.columbia.edu/tag/0DPA}.

\bibitem{Varolin}
D.~Varolin, Division theorems and twisted complexes,
\emph{Math. Z.} \textbf{259} (2008), no.~1, 1--20.
\href{https://doi.org/10.1007/s00209-007-0133-4}{doi:10.1007/s00209-007-0133-4}.

\end{thebibliography}
\end{document}
